% Self-contained body. Requires amsmath, amssymb, amsthm and theorem environments.
% This fragment is written for integration by the parent agent.

\section{The complete critical law for linear recovery}

Fix a prime power $q$ and let $n\ge4$ be even. Put
\[
 V=\mathbb F_q^n,\qquad E=\operatorname{Alt}(V),\qquad
 D=\binom n2,\qquad d=D-1,
\]
where alternating means $G(v,v)=0$ for all $v$, also in characteristic two.
Fix a nondegenerate alternating form $H$. Let $Q_{n,t}$ be the law of
$t$ independent uniform two-dimensional subspaces of $V$, and let
$P_{H,n,t}$ be the law of $t$ independent uniform $H$-isotropic
two-dimensional subspaces. Define
\[
 N_t=\{G\in E:G|_{K_i\times K_i}=0\ (1\le i\le t)\}.
\]
Under $P_H$, the line $\langle H\rangle$ is contained in $N_t$.
The mixture $P_{n,t}$ chooses $[H]$ uniformly from the nondegenerate
projective forms and then samples according to $P_{H,n,t}$.

For $j\ge0$ and $b\in\mathbb Z$, write
\[
 g_j=\prod_{i=1}^j(1-q^{-i}),\qquad
 g_\infty=\prod_{i=1}^{\infty}(1-q^{-i}),
\]
and define the probability mass function
\begin{equation}\label{crit:pi}
 \pi_b(k)=\frac{g_\infty q^{-k(k+b)}}{g_k g_{k+b}},
 \qquad k\ge\max(0,-b),
\end{equation}
with value zero outside this range.

\begin{theorem}[Critical nullity universality]\label{crit:main}
For each fixed integer $a$, as even $n\to\infty$,
\begin{align}
 \mathcal L_{Q_{n,d+a}}(\dim N_{d+a})&\longrightarrow\pi_{a-1},
 \label{crit:null}\\
 \mathcal L_{P_{H,n,d+a}}(\dim N_{d+a}-1)&\longrightarrow\pi_a
 \label{crit:plant}
\end{align}
in total variation. The second convergence is uniform over the
nondegenerate form $H$, and therefore holds under the mixture $P$.
Let
\[
 \rho_q=\frac{2q-1}{q^2},\qquad
 \beta_q=-\log_q\rho_q>0.
\]
For fixed $q,a$, both errors are
\begin{equation}\label{crit:rate}
 O_{q,a}\!\left(n^2q^{-\beta_q n/8}\right).
\end{equation}
More generally, the joint nullity process over any fixed finite set of
integer offsets converges to the corresponding uniform-matrix rank
process.
\end{theorem}

The mass function \eqref{crit:pi} is the classical rectangular uniform
finite-field matrix law~\cite{FulmanGoldstein2015}. The assertion here is its universality for
rows constrained to be decomposable bivectors, and for their further
restriction to a symplectic hyperplane. These rows have strongly
dependent coordinates and are not uniform vectors in the ambient
space.

\begin{corollary}[Exact linear-recovery curve]\label{crit:recovery}
At $t=d+a$, ordinary Gaussian elimination uniquely determines the
hidden projective form, in the sense $N_t=\langle H\rangle$, with
limiting probability
\begin{equation}\label{crit:recoverycurve}
 F_q(a)=
 \begin{cases}
 0,&a<0,\\
 \displaystyle\frac{g_\infty}{g_a}
   =\prod_{j=a+1}^{\infty}(1-q^{-j}),&a\ge0.
 \end{cases}
\end{equation}
The error for fixed $a$ satisfies \eqref{crit:rate}.
\end{corollary}

The event in this corollary is a certificate supplied by linear algebra.
It need not coincide with unique recovery among nondegenerate forms:
a higher-dimensional $N_t$ could contain just one nondegenerate
projective form. No such equivalence is used.

\begin{corollary}[Exact dimension-only testing curve]\label{crit:testing}
Let $\Delta_{n,a}$ be the total variation distance between the laws of
$\dim N_{d+a}$ under $P_{n,d+a}$ and $Q_{n,d+a}$. Then
\begin{equation}\label{crit:testcurve}
 \Delta_{n,a}\longrightarrow
 T_q(a)=
 \begin{cases}
 \displaystyle q^{a-1}\frac{g_\infty}{g_{1-a}},&a\le0,\\[2mm]
 \displaystyle\frac{g_\infty}{g_{a-1}},&a\ge1.
 \end{cases}
\end{equation}
A test attaining this limiting advantage declares the planted model
when $\dim N\ge2-a$ for $a\le0$, and when $\dim N\ge1$ for $a\ge1$.
These claims concern the dimension statistic; they do not identify the
optimal test using all observations.
\end{corollary}

\begin{corollary}[The observation overhead]\label{crit:overhead}
Sample under a fixed hidden form until the first time $T_n$ at which
$N_{T_n}=\langle H\rangle$. Then $T_n-d$ converges in distribution,
and in every fixed moment, to a nonnegative integer $C_q$ with
\[
 \Pr(C_q=a)=\frac{g_\infty q^{-a}}{g_a},\qquad a\ge0.
\]
Equivalently,
\[
 C_q\ \stackrel{\mathrm d}=\ \sum_{j=1}^{\infty}F_j,
 \qquad
 \Pr(F_j=r)=(1-q^{-j})q^{-jr}\quad(r\ge0),
\]
where the $F_j$ are independent. The sum is almost surely finite, and
\begin{align*}
 \lim_{n\to\infty}\bigl(\mathbb ET_n-d\bigr)
 &=\sum_{j\ge1}\frac1{q^j-1},\\
 \lim_{n\to\infty}\operatorname{Var}(T_n)
 &=\sum_{j\ge1}\frac{q^j}{(q^j-1)^2}.
\end{align*}
\end{corollary}

\section{Fourier coefficients of the actual observations}

Choose a uniformly random ordered basis $(u,v)$ of each observed
two-plane. This extra choice does not change any kernel $N_t$.
Under $Q$, $(u,v)$ is uniform among independent ordered pairs; under
$P_H$, it is uniform among independent ordered pairs satisfying
$H(u,v)=0$. The row defined by this pair is the linear functional
$G\mapsto G(u,v)$ on $E$, or on $E/\langle H\rangle$ under $P_H$.

Fix a nontrivial additive character $\psi$ of $\mathbb F_q$. Let
\[
 \delta_n=(q+1)q^{-n}-q^{1-2n}
\]
be the probability that two independent uniform vectors of $V$ are
linearly dependent. Unconditionally,
\begin{equation}\label{crit:basicFourier}
 \mathbb E_{u,v}\psi(G(u,v))=q^{-\operatorname{rank}G}.
\end{equation}
Indeed, averaging over $v$ gives one if $u\in\ker G$ and zero otherwise.
For dependent pairs every alternating form vanishes. Consequently,
under the independent-pair null law,
\begin{equation}\label{crit:nullFourier}
 \phi_Q(G)=
 \frac{q^{-\operatorname{rank}G}-\delta_n}{1-\delta_n}.
\end{equation}
Additive-character orthogonality gives the planted formula
\begin{equation}\label{crit:plantFourier}
 \phi_H(G)=
 \frac{\displaystyle\sum_{c\in\mathbb F_q}
 q^{-\operatorname{rank}(G+cH)}-q\delta_n}
 {(1-q^{-n})(1-q^{2-n})}.
\end{equation}
To verify the denominator, before excluding dependent pairs,
\[
 \Pr\{H(u,v)=0\}=q^{-1}\bigl(1+(q-1)q^{-n}\bigr),
\]
so after exclusion, multiplication by $q$ gives exactly
$(1-q^{-n})(1-q^{2-n})$.

In the null model put $W=E$, $m=D$, and $r(G)=\operatorname{rank}G$.
In the planted model put $W=E/\langle H\rangle$, $m=d$, and
\[
 r(\overline G)=\min_{c\in\mathbb F_q}\operatorname{rank}(G+cH).
\]
This is well defined on the quotient and is invariant under nonzero
scalar multiplication. For a nonzero $x\in W$, $r(x)\ge2$.
Equations \eqref{crit:nullFourier}--\eqref{crit:plantFourier} imply, for
$n\ge4$, $R\le n$, and $r(x)>R$,
\begin{equation}\label{crit:FourierBound}
 |\mathbb E\psi(\ell(x))|\le C_q q^{-R},\qquad
 C_q=2q(q+2).
\end{equation}
The same estimate holds for every nontrivial additive character, since
replacing $\psi$ by $z\mapsto\psi(sz)$ replaces $G$ by $sG$ and does
not change $r(x)$. To check the constant, both displayed denominators
are at least $1/2$. In the planted numerator use
$\sum_cq^{-\operatorname{rank}(G+cH)}\le q q^{-R}$ and
$q\delta_n\le q(q+1)q^{-n}\le q(q+1)q^{-R}$. The null numerator is
smaller under the same bound.

\begin{lemma}[A uniform survival bound]\label{crit:rho}
For every nonzero $x\in W$ in either model,
\[
 h(x):=\Pr\{\ell(x)=0\}\le\rho_q=\frac{2q-1}{q^2}<1.
\]
If $r(x)>R$, then
\begin{equation}\label{crit:survivalFine}
 h(x)\le q^{-1}\exp(K_q q^{-R}),\qquad
 K_q=(q-1)C_q.
\end{equation}
\end{lemma}

\begin{proof}
In the null model \eqref{crit:nullFourier} and character orthogonality
give
\[
 h(x)=q^{-1}+(1-q^{-1})
 \frac{q^{-r(x)}-\delta_n}{1-\delta_n}
 \le q^{-1}+(1-q^{-1})q^{-2}\le\rho_q.
\]
In the planted model fix a representative $G\notin\langle H\rangle$.
For a nonzero vector $u$, the functionals $G(u,\cdot)$ and
$H(u,\cdot)$ are dependent exactly when
$(G-cH)(u,\cdot)=0$ for one, necessarily unique, scalar $c$.
Every $G-cH$ is a nonzero alternating form, hence has rank at least two.
The fraction $\theta$ of such nonzero $u$ therefore satisfies
\[
 \theta=\frac{\sum_c(q^{\dim\ker(G-cH)}-1)}{q^n-1}
 \le\frac{q(q^{n-2}-1)}{q^n-1}<q^{-1}.
\]
Conditional on a fixed nonzero $u$ and on $H(u,v)=0$, independence
leaves $q^{n-1}-q$ choices for $v$. In the dependent-functional case
all are $G$-isotropic; otherwise precisely $q^{n-2}-q$ are.
Thus
\[
 h(x)=\theta+(1-\theta)\frac{q^{n-2}-q}{q^{n-1}-q}
 \le q^{-1}+(1-q^{-1})q^{-1}=\rho_q.
\]
Finally, character orthogonality and \eqref{crit:FourierBound} give
$h(x)\le q^{-1}(1+(q-1)C_q q^{-R})$, proving
\eqref{crit:survivalFine}.
\end{proof}

\section{The burn-in argument and the terminal coupling}

The number $B_{n,R}$ of alternating forms with rank at most $R$ obeys
\begin{equation}\label{crit:rankcount}
 B_{n,R}\le g_\infty^{-1}q^{D-\binom{n-R}{2}}.
\end{equation}
Indeed, such a form has a kernel containing an $(n-R)$-space. The
number of these spaces is at most
$g_\infty^{-1}q^{R(n-R)}$, and a form vanishing there descends to
one of $q^{R(R-1)/2}$ alternating forms on the quotient.
In the planted model the number of nonzero quotient vectors with
$r(x)\le R$ is at most $B_{n,R}$: each is the image of at least one
rank-at-most-$R$ form. The same bound is immediate in the null model.

Let $1\le L\le m$, $0\le R<n$, and put $t_0=m-L$. Let
$A_s\subseteq W$ denote the common kernel after $s$ observations, so
that $A_s=N_s$ in the null model and $A_s=N_s/\langle H\rangle$ in
the planted model. Necessarily $\dim A_{t_0}\ge L$. By the preceding
survival bounds,
\begin{align}
 \mathbb E q^{\dim A_{t_0}}
 &=1+\sum_{x\ne0}h(x)^{t_0}\nonumber\\
 &\le1+q^L e^{\eta}+B_{n,R}\rho_q^{t_0},\qquad
 \eta=K_qt_0q^{-R}.\label{crit:firstmoment}
\end{align}
Since $q^k-q^L\ge(q-1)q^L$ for $k>L$, this proves
\begin{equation}\label{crit:burnRank}
 \Pr\{\dim A_{t_0}>L\}
 \le\frac{q^{-L}+e^\eta-1+q^{-L}B_{n,R}\rho_q^{t_0}}{q-1}.
\end{equation}
A separate union bound proves
\[
 \Pr\{\exists x\in A_{t_0}\setminus\{0\}:r(x)\le R\}
 \le B_{n,R}\rho_q^{t_0}.
\]
Hence the event $\mathcal E$ that $\dim A_{t_0}=L$ and every nonzero
element of $A_{t_0}$ has $r(x)>R$ has failure probability at most
\begin{equation}\label{crit:burnError}
 e_{n,L,R}:=
 \frac{q^{-L}+e^\eta-1}{q-1}+2B_{n,R}\rho_q^{t_0}.
\end{equation}

On $\mathcal E$, every later kernel $A$ inherits this rank-distance
property. If $\dim A=k$, Fourier orthogonality on the vector space $A$
gives
\[
 \Pr\{\ell|_A=0\}
 =q^{-k}\sum_{x\in A}\mathbb E\psi(\ell(x)).
\]
Its difference from $q^{-k}$ is at most $C_qq^{-R}$, by
\eqref{crit:FourierBound}. This bound has no factor growing with $k$.
The next kernel dimension is $k$ on this event and $k-1$ otherwise.
For a uniform row in $\mathbb F_q^L$, precisely the same transition
has probability $q^{-k}$ of staying at $k$. Coupling the two Bernoulli
transitions successively, for every integer $b$ with $L+b\ge0$,
\begin{equation}\label{crit:coupling}
 d_{\mathrm{TV}}\!\left(
 \mathcal L(\dim A_{m+b}),\nu_{L,b}\right)
 \le e_{n,L,R}+(L+b)C_qq^{-R},
\end{equation}
where $\nu_{L,b}$ is the nullity law of an $(L+b)\times L$ matrix
of independent uniform entries. Conditioning on $\mathcal E$ causes
no problem: future observations are independent of the burn-in.
The same coupling controls the whole terminal dimension path.

\subsection{The uniform-matrix limit, including its error}

Counting rank-$L-k$ matrices gives
\[
 \nu_{L,b}(k)=q^{-k(k+b)}
 \frac{g_Lg_{L+b}}{g_kg_{k+b}g_{L-k}}.
\]
The allowed range is $\max(0,-b)\le k\le L$, and the count is zero
outside it. One derivation counts a full-rank $(L+b)\times(L-k)$
matrix and a full-rank $(L-k)\times L$ matrix, then divides by
$|\operatorname{GL}_{L-k}(q)|$ for their common change of basis.
The displayed probabilities converge pointwise to \eqref{crit:pi}.
They are tight because
\[
 \mathbb E q^{\operatorname{nullity}}=1+(q^L-1)q^{-(L+b)}
 \le1+q^{-b}.
\]
Thus \eqref{crit:pi} sums to one and the convergence is in total
variation.

An elementary coupling supplies a useful uniform error. In a uniform
$(L+h+b)\times(L+h)$ matrix, expose its first $h$ rows. The probability
that they are not independent is at most
\[
 \sum_{j=0}^{h-1}q^{j-(L+h)}\le\frac{q^{-L}}{q-1}.
\]
When they are independent, the remaining $L+b$ rows restrict to
independent uniform linear functionals on an $L$-dimensional kernel.
Therefore
\[
 d_{\mathrm{TV}}(\nu_{L+h,b},\nu_{L,b})\le\frac{q^{-L}}{q-1}.
\]
Letting $h\to\infty$ gives
\begin{equation}\label{crit:uniformError}
 d_{\mathrm{TV}}(\nu_{L,b},\pi_b)\le\frac{q^{-L}}{q-1}.
\end{equation}

\subsection{An explicit finite error bound}

Combining \eqref{crit:coupling} and \eqref{crit:uniformError} proves
\begin{equation}\label{crit:finite}
 d_{\mathrm{TV}}\!\left(\mathcal L(\dim A_{m+b}),\pi_b\right)
 \le
 \frac{2q^{-L}+e^{K_q(m-L)q^{-R}}-1}{q-1}
 +2B_{n,R}\rho_q^{m-L}+(L+b)C_qq^{-R}.
\end{equation}
As always, the right side may be capped at one. This holds separately
for each admissible $L,R$; they may be optimized for any finite $n$.

For the asymptotic assertion take
$L=R=\lfloor\beta_q n/8\rfloor$. For sufficiently large $n$ these
parameters are admissible. Since $m\ge D-1$,
\begin{align*}
 \log_q\bigl(g_\infty B_{n,R}\rho_q^{m-L}\bigr)
 &\le D-\binom{n-R}{2}-\beta_q(m-L)\\
 &\le nR-\beta_q(D-1-L)\\
 &\le-\beta_q n^2/4
\end{align*}
for $n\ge8$, using $0<\beta_q<1$ and $L=R\le\beta_q n/8$.
The remaining terms in \eqref{crit:finite} are
$O_{q,b}(n^2q^{-\beta_q n/8})$. Choosing $m=D$, $b=a-1$ proves
\eqref{crit:null}; choosing $m=d$, $b=a$ proves
\eqref{crit:plant}. All constants are independent of $H$.

\section{Proofs of the statistical consequences}

Equation \eqref{crit:recoverycurve} is \eqref{crit:plant} evaluated at
zero. For $a<0$, linear dimension alone makes this event impossible.

For the dimension test, let $k=\dim N$. On the support of the null
limit law,
\begin{equation}\label{crit:dimensionLikelihood}
 \frac{\pi_a(k-1)}{\pi_{a-1}(k)}
 =q^{a-1}(q^k-1),
\end{equation}
where the numerator is interpreted as zero at $k=0$. This follows
immediately from $g_k/g_{k-1}=1-q^{-k}$.
If $a\ge1$, this ratio is at least one for every $k\ge1$ and is
zero at $k=0$. Total variation is consequently
$\pi_{a-1}(0)=g_\infty/g_{a-1}$.
If $a\le0$, the minimum possible dimension is $k_0=1-a$.
At $k_0$ the ratio is $1-q^{a-1}<1$ and it exceeds one at every
larger integer. Thus total variation equals
$q^{a-1}\pi_{a-1}(1-a)=q^{a-1}g_\infty/g_{1-a}$.
Total variation convergence in the main theorem justifies taking
these limiting distances and the proposed tests.

Finally, $\Pr(T_n-d\le a)=\Pr(N_{d+a}=\langle H\rangle)$ for $a\ge0$,
so the limiting distribution function is $g_\infty/g_a$.
Taking differences gives the mass function in
Corollary~\ref{crit:overhead}, including the value at zero.
For an ordinary uniform matrix, the number of rows spent at kernel
dimension $j$ before the next dimension reduction is geometric with
success probability $1-q^{-j}$, independently for successive $j$.
Subtracting the one indispensable row at each dimension gives the
independent $F_j$ representation. Since
$\sum_j\Pr(F_j>0)=\sum_jq^{-j}<\infty$, the infinite sum is almost
surely finite; its mean and variance are the displayed convergent
series.

It remains to justify moment convergence for the actual observation
model. A union bound over projective points of $W$ gives, at $t=d+a$,
\[
 \Pr_H(T_n-d>a)
 \le\frac{q^{-a}}{q-1}\exp\bigl(K_q(d+a)q^{-R}\bigr)
   +\frac{B_{n,R}}{q-1}\rho_q^{d+a}.
\]
For $R=\lfloor\beta_qn/8\rfloor$ and all sufficiently large $n$,
$K_qq^{-R}\le(\log q)/2$, $K_qdq^{-R}\le1$, and
$B_{n,R}\rho_q^d\le1$. Hence uniformly for all $a\ge0$,
\[
 \Pr_H(T_n-d>a)
 \le\frac{e}{q-1}q^{-a/2}+\frac1{q-1}\rho_q^a.
\]
This summable exponential bound implies uniform integrability of every
fixed power, completing the proof.

\paragraph{What remains open.}
The theorem determines the critical nullity process, the exact success
curve for its linear certificate, and the complete testing curve after
retaining only the dimension. It does not determine the number of
nondegenerate projective forms inside $N_t$ or the full likelihood
distribution. Those depend on the restriction of the Pfaffian to a
small random kernel and remain distinct questions.
